\section{Elección tecnológica}
\subsection{Sistema de archivos o base de datos}

Para el prototipo solicitado resulta conveniente utilizar un sistema de
archivos. La información se limita a tres entidades, el acceso se realiza desde
una aplicación de consola y el objetivo inmediato es disponer de persistencia
sin esperar el diseño completo de la base de datos. Los CSV también facilitan
la inspección de los registros y su posterior importación, siempre que se
conserven encabezados, codificación y reglas de escape consistentes
\cite{shafranovich2005}.

Para la operación definitiva de PuellaGame, la alternativa más conveniente es
una base de datos administrada por un SMBD. El crecimiento de clientes,
sucursales y premios exigirá relaciones entre entidades, restricciones de
integridad, consultas más variadas, control de acceso, respaldo y trabajo
simultáneo de varios usuarios. Estos servicios forman parte del propósito de un
gestor y evitan que cada aplicación tenga que implementarlos de manera aislada
\cite{elmasri2016,silberschatz2019}.

Por ello, la decisión depende del horizonte de uso: los archivos CSV son una
solución transitoria adecuada para el avance urgente; la base de datos es la
solución sostenible para el sistema real. El prototipo reduce el costo de la
transición al conservar una llave estable por entidad, validar los datos desde
su captura y separar la persistencia del resto de la aplicación.
